\section{Descrição das unidades}
\label{sec:unidades}

\subsection{ALU}
\label{sec:alu}

A ALU (\sig{alu.vhd}) é um circuito puramente combinacional: não recebe clock
e não guarda estado. Suas entradas são os operandos \sig{ALU\_A} e
\sig{ALU\_B} e o comando \sig{CMD}; suas saídas são o resultado \sig{ALU\_R} e o
vetor de seis flags. A UC só envia operandos à ALU nas instruções lógicas e
aritméticas; nas demais, as entradas são mantidas em zero. Cabe à UC salvar o
resultado e as flags no fim do estado de \emph{execute}.

A codificação de \sig{CMD} foi definida no pacote \sig{cpu\_pkg}
(Tabela~\ref{tab:alucmd}). Uma única ALU atende a todas as instruções: \instr{inc}
e \instr{dec} usam os comandos de soma e subtração com \sig{ALU\_B}~=~1, e as
operações de um operando usam \sig{ALU\_B}~=~0.

\begin{table}[H]
\centering
\caption{Comandos da ALU e instruções que os usam.}
\label{tab:alucmd}
\small
\begin{tabular}{l l l l l}
\toprule
\sig{CMD} & Operação & Resultado em 9 bits ($r_8 \ldots r_0$) & \sig{ALU\_B} & Instruções \\
\midrule
\texttt{0000} & soma        & $A + B$                       & \sig{Ry} / 1 & \instr{add}, \instr{inc} \\
\texttt{0001} & subtração   & $A - B$                       & \sig{Ry} / 1 & \instr{sub}, \instr{dec} \\
\texttt{0010} & E lógico    & $0\,(A \land B)$              & \sig{Ry} & \instr{and} \\
\texttt{0011} & OU lógico   & $0\,(A \lor B)$               & \sig{Ry} & \instr{or} \\
\texttt{0100} & OU exclusivo& $0\,(A \oplus B)$             & \sig{Ry} & \instr{xor} \\
\texttt{0101} & negação     & $0\,\lnot A$                  & 0 & \instr{not} \\
\texttt{0110} & rotação esq.& $A_7\,A_{6..0}\,A_7$          & 0 & \instr{rol} \\
\texttt{0111} & rotação dir.& $A_0\,A_0\,A_{7..1}$          & 0 & \instr{ror} \\
\texttt{1000} & desl. esq.  & $A_7\,A_{6..0}\,0$            & 0 & \instr{lsl} \\
\texttt{1001} & desl. dir.  & $A_0\,0\,A_{7..1}$            & 0 & \instr{lsr} \\
\bottomrule
\end{tabular}
\end{table}

\paragraph{Conta em 9 bits.} Os operandos são estendidos com um zero à
esquerda (interpretação \emph{unsigned}) e a operação é feita em 9 bits. Os 8
bits de baixo formam o resultado e o bit 8 é diretamente a flag OVERFLOW: na
soma é o \emph{carry}; na subtração é o \emph{borrow} (1 quando $A < B$). Nas
rotações e deslocamentos, o bit 8 é preenchido com o bit que ``sai'' do
registrador, e nas operações lógicas, com 0. Assim, uma única expressão
(\sig{OVERFLOW <= r9(8)}) atende a todas as convenções do enunciado.

\begin{lstlisting}[style=vhdl, caption={Núcleo da ALU (trecho de \texttt{alu.vhd}).}]
a9 <= '0' & unsigned(ALU_A);   -- estende com zero: 0..255
b9 <= '0' & unsigned(ALU_B);
...
case CMD is
    when ALU_ADD => r9 <= a9 + b9;                    -- add, inc
    when ALU_SUB => r9 <= a9 - b9;                    -- sub, dec
    when ALU_ROL => r9 <= unsigned(ALU_A(7) & ALU_A(6 downto 0) & ALU_A(7));
    when ALU_LSR => r9 <= unsigned(ALU_A(0) & '0' & ALU_A(7 downto 1));
    ...
end case;
ALU_R         <= std_logic_vector(r9(7 downto 0));
FLAGS(F_ZERO) <= '1' when r9(7 downto 0) = 0 else '0';
FLAGS(F_NEG)  <= r9(7);
FLAGS(F_OVF)  <= r9(8);
FLAGS(F_EQ)   <= '1' when unsigned(ALU_A) = unsigned(ALU_B) else '0';
FLAGS(F_GT)   <= '1' when unsigned(ALU_A) > unsigned(ALU_B) else '0';
FLAGS(F_SM)   <= '1' when unsigned(ALU_A) < unsigned(ALU_B) else '0';
\end{lstlisting}

\paragraph{Flags.} A Tabela~\ref{tab:flags} resume o significado de cada flag.
OVERFLOW, GREATER e SMALLER interpretam os operandos como \emph{unsigned};
NEGATIVE é o bit de sinal do resultado em complemento de 2. Como soma e
subtração produzem os mesmos 8 bits para valores com ou sem sinal, cabe ao
programador escolher a flag adequada: para comparar valores sem sinal usa-se
\instr{bcs}/\instr{bcc} ou \instr{blt}/\instr{bgt}, e não \instr{bneg} --- por
exemplo, $200 - 10 = 190$ tem NEGATIVE = 1, embora $200 > 10$.

\begin{table}[H]
\centering
\caption{Flags da ALU (\emph{A} e \emph{B} são \sig{ALU\_A} e \sig{ALU\_B}).}
\label{tab:flags}
\small
\begin{tabularx}{\linewidth}{l l X}
\toprule
Flag & Definição & Observação \\
\midrule
ZERO     & resultado $= 0$            & \\
NEGATIVE & resultado$[7] = 1$         & bit de sinal em complemento de 2 \\
OVERFLOW & bit 8 da conta em 9 bits    & \emph{carry} em \instr{add}/\instr{inc}; \emph{borrow} em \instr{sub}/\instr{dec}; bit que sai em \instr{rol}/\instr{lsl} ($A_7$) e \instr{ror}/\instr{lsr} ($A_0$); 0 nas lógicas \\
EQUAL    & $A = B$                    & em \instr{inc}/\instr{dec} compara com 1; nas de um operando, com 0 \\
GREATER  & $A > B$ (sem sinal)        & idem \\
SMALLER  & $A < B$ (sem sinal)        & idem \\
\bottomrule
\end{tabularx}
\end{table}

Exemplos: $3 - 5$ em 9 bits resulta $510 = 1\,1111\,1110_2$, ou seja, resultado
\texttt{0xFE} com OVERFLOW = 1; $200 + 100 = 300 = 1\,0010\,1100_2$ resulta 44
com OVERFLOW = 1; \instr{inc} de 255 resulta 0 com OVERFLOW = 1 e ZERO = 1. A
conta não usa o tipo \sig{signed}: a extensão de sinal faria o bit 8 deixar de
ser o \emph{carry}.

\subsection{Unidade de Controle}
\label{sec:uc}

A Unidade de Controle está em \sig{cpu.vhd}, num único processo síncrono que
implementa a máquina de estados da Seção~\ref{sec:fsm}. Ela tem três
responsabilidades:

\begin{enumerate}[nosep]
  \item \textbf{Decodificar a instrução.} O byte do \sig{IR} é convertido
        combinacionalmente no tipo enumerado \sig{instr\_t} pela função
        \sig{decode}, e os campos \sig{rx} e \sig{ry} são extraídos dos bits
        [3:2] e [1:0]. Todas as decisões da UC são tomadas sobre esse tipo, o
        que torna o código legível (\sig{when I\_PUSH =>} \ldots) e garante que
        códigos inválidos caiam em \sig{I\_NOP}.
  \item \textbf{Gerar os sinais de controle}: seleção dos multiplexadores de
        \sig{MAR}, \sig{MBR} e \sig{PC}, operandos e comando da ALU,
        habilitação de escrita na RAM (\sig{WE}) e escrita nos registradores.
  \item \textbf{Guardar as flags.} As seis flags são copiadas da ALU para o
        registrador \sig{FLAGS} (flip-flops dentro da UC) no fim do
        \emph{execute}, e apenas nas instruções de ALU. Elas permanecem
        inalteradas durante instruções de memória, saltos, \instr{mov},
        \instr{nop} e \instr{halt}, e são lidas pelos saltos condicionais e
        exibidas nos LEDs.
\end{enumerate}

O processo é sensível apenas ao clock de \SI{50}{\mega\hertz}; o reset é
síncrono e tem prioridade, e fora dele a FSM só avança quando \sig{CE}~=~1:

\begin{lstlisting}[style=vhdl, caption={Estrutura do processo da UC (trecho de \texttt{cpu.vhd}).}]
p_uc : process(CLK)
begin
    if rising_edge(CLK) then
        if RESET = '1' then
            PC <= (others => '0');  SP <= to_unsigned(254, 8);
            FLAGS <= (others => '0');  state <= S_FETCH;  ...
        elsif CE = '1' then
            case state is
                when S_FETCH   => ...  state <= S_DECODE;
                when S_DECODE  => ...  state <= S_EXECUTE;
                when S_EXECUTE => ...  state <= S_WRITE;
                when S_WRITE   => ...  state <= S_FETCH;
            end case;
        end if;
    end if;
end process;
\end{lstlisting}

\subsection{Registradores}

Todos os registradores são sinais de 8 bits atualizados apenas dentro do
processo da UC, e portanto sintetizados como flip-flops com \emph{enable}. O
banco de registradores é um vetor de quatro palavras (\sig{REG(0 to 3)}),
indexado por \sig{rx} e \sig{ry}; a leitura é combinacional (dois
multiplexadores 4:1) e a escrita ocorre apenas no estado WRITE, sempre a partir
de \sig{MBR}.

\begin{itemize}[nosep]
  \item \sig{PC} é do tipo \sig{unsigned}, o que permite somar 1 ou 2
        diretamente; o estouro de 255 para 0 é natural.
  \item \sig{IR} recebe o byte da instrução no FETCH. No reset ele é iniciado
        com \texttt{0xF0} (\instr{nop}) para que o LCD não mostre uma instrução
        inexistente antes da primeira busca.
  \item \sig{MAR} e \sig{MBR} são as únicas ligações com a RAM
        (\sig{RAM\_ADDR <= MAR}, \sig{RAM\_DIN <= MBR}).
  \item \sig{WE} também é um registrador: é ligado no DECODE das instruções que
        escrevem na memória e desligado no EXECUTE.
\end{itemize}

\subsection{Memória RAM}
\label{sec:ram}

A RAM (\sig{RAM.vhd}) mantém a interface fornecida (\sig{CLK}, \sig{DIN},
\sig{ADDR}, \sig{WE}, \sig{DOUT}, \sig{POS\_255}) e é inferida como
\emph{block RAM}. Ela é síncrona: na borda de \textbf{descida} do clock ela
grava \sig{DIN} se \sig{WE}~=~1 e registra o endereço de leitura; \sig{DOUT} é o
conteúdo desse endereço registrado. A saída \sig{POS\_255} expõe
permanentemente a posição 255, usada como saída para o LCD.

O programa é o valor inicial do vetor \sig{ram}, no formato
\sig{endereço => "bbbbbbbb"}. Para executar outro programa (por exemplo, o
programa surpresa) basta substituir esse bloco e sintetizar novamente; a
interface da entidade não muda.

\paragraph{Latência de leitura.} Como o endereço só é registrado na descida do
clock, o dado não está disponível no mesmo instante em que o endereço é
aplicado. A UC trata isso com uma regra simples: \emph{todo endereço é colocado
em \sig{MAR} num estado, e o dado é lido de \sig{DOUT} no estado seguinte}
(Figura~\ref{fig:ramtiming}). Como a UC muda \sig{MAR} na borda de subida, a
borda de descida seguinte (\SI{10}{\nano\second} depois) já registra o novo
endereço, e o dado fica válido muito antes da próxima borda em que \sig{CE}~=~1.
A regra continua válida mesmo sem divisor de clock (\sig{CE} sempre 1), porque
há sempre uma borda de descida entre duas bordas de subida.

\begin{figure}[H]
\centering
\begin{tikzpicture}[x=1.1cm, y=0.55cm, thick]
  \def\T{2}
  % clock
  \node[anchor=east, font=\scriptsize\ttfamily] at (-0.2,6.5) {CLK};
  \draw (0,6) \foreach \i in {0,...,5} { -- ++(0,1) -- ++(1,0) -- ++(0,-1) -- ++(1,0) };
  % CE (sempre 1 neste exemplo)
  \node[anchor=east, font=\scriptsize\ttfamily] at (-0.2,4.75) {estado};
  \draw (0,4.25) -- (12,4.25) (0,5.25) -- (12,5.25);
  \foreach \x/\t in {0/FETCH,2/DECODE,4/EXECUTE,6/WRITE,8/FETCH,10/DECODE} {
     \draw (\x,4.25) -- (\x,5.25); \node[font=\tiny] at (\x+1,4.75) {\t}; }
  % MAR
  \node[anchor=east, font=\scriptsize\ttfamily] at (-0.2,3.0) {MAR};
  \draw (0,2.5) -- (12,2.5) (0,3.5) -- (12,3.5);
  \foreach \x/\t in {0/PC,2/PC+1,4/Ry (endereço),8/PC',10/PC'+1} {
     \draw (\x,2.5) -- (\x,3.5); \node[font=\tiny] at (\x+1,3.0) {\t}; }
  % DOUT
  \node[anchor=east, font=\scriptsize\ttfamily] at (-0.2,1.25) {DOUT};
  \draw (0,0.75) -- (12,0.75) (0,1.75) -- (12,1.75);
  \foreach \x/\t in {1/MEM[PC],3/MEM[PC+1],5/MEM[Ry],9/MEM[PC']} {
     \draw (\x,0.75) -- (\x,1.75); \node[font=\tiny] at (\x+1,1.25) {\t}; }
  \draw (11,0.75) -- (11,1.75);
  % setas
  \draw[-{Latex}, red!70!black, thin] (2,0.6) -- (2,-0.3) node[below, font=\tiny, align=center] {IR $\gets$ DOUT};
  \draw[-{Latex}, red!70!black, thin] (4,0.6) -- (4,-1.9) node[below, font=\tiny, align=center] {MBR $\gets$ DOUT\\(se \instr{ldi}/\instr{jmp})};
  \draw[-{Latex}, red!70!black, thin] (6,0.6) -- (6,-0.3) node[below, font=\tiny, align=center] {MBR $\gets$ DOUT\\(\instr{ldr}, \instr{pop})};
\end{tikzpicture}
\caption{Temporização das leituras da RAM síncrona durante um \instr{ldr Rx, [Ry]} (exemplo com \sig{CE}
sempre 1; com o divisor, cada estado dura \SI{0,5}{\second}, mas a relação entre
as bordas é a mesma). \sig{DOUT} muda sempre na borda de descida do clock.}
\label{fig:ramtiming}
\end{figure}

\paragraph{Escrita.} Nas instruções \instr{st}, \instr{str} e \instr{push}, a UC
coloca o endereço em \sig{MAR}, o dado em \sig{MBR} e liga \sig{WE} no DECODE;
\sig{WE} fica em 1 durante todo o EXECUTE e é desligado na borda que encerra esse
estado. Com o divisor ativo, a RAM vê várias bordas de descida com
\sig{WE}~=~1, mas sempre com o mesmo endereço e o mesmo dado, de modo que as
escritas repetidas são idempotentes.

\subsection{Divisor de clock (\emph{clock enable})}
\label{sec:ce}

O enunciado pede que cada instrução completa leve cerca de \SI{2}{\second}. Como
todas as instruções passam pelos mesmos quatro estados, cada estado deve durar
\SI{0,5}{\second}, ou seja,
\[
  \text{TICKS} = \SI{0,5}{\second} \times \SI{50}{\mega\hertz} = \num{25000000}
  \text{ ciclos por estado.}
\]
O módulo \sig{clock\_enable} conta de 0 a $\text{TICKS}-1$ e gera
\sig{CE}~=~1 durante um único ciclo no fim da contagem. O valor é um
\emph{generic}, o que permite usar valores pequenos na simulação. O mesmo
contador gera o sinal do LED CLOCK: alto na primeira metade de cada estado e
baixo na segunda, de modo que cada pulso do LED corresponde a uma transição da
UC.

Usar um \emph{clock enable} em vez de dividir o clock tem duas vantagens: a
CPU continua no mesmo domínio de clock da RAM e do LCD, e não se cria um clock
gerado por lógica, que a FPGA roteia mal e que causaria problemas de
temporização.

\subsection{Stack pointer}

O registrador \sig{SP} aponta para a próxima posição livre da pilha, que cresce
para endereços menores. Ele é iniciado em 254 no reset, dentro da UC, porque a
posição 255 é reservada para saída.

\begin{itemize}[nosep]
  \item \instr{push Rx}: no DECODE, $MAR \gets SP$, $MBR \gets Rx$ e
        $WE \gets 1$; a escrita ocorre no EXECUTE; no WRITE, $SP \gets SP - 1$.
  \item \instr{pop Rx}: no DECODE, $MAR \gets SP + 1$; no EXECUTE,
        $MBR \gets DOUT$; no WRITE, $Rx \gets MBR$ e $SP \gets SP + 1$.
\end{itemize}

Assim, após \instr{push Ra}; \instr{push Rb} a partir do reset, $MEM[254] = Ra$,
$MEM[253] = Rb$ e $SP = 252$; dois \instr{pop} devolvem os valores na ordem
inversa e restauram $SP = 254$. Esse comportamento é verificado pelo teste T8
(Seção~\ref{sec:programas}).
